\chapter{范畴论：保留结构的箭头与谱投影}
\label{chap:category}

前面一直在做一种比较：两种写法是否描述了同一个几何结构？范畴论把比较所允许的变换称为\textbf{箭头}。先规定对象和允许的箭头，再问哪些对象同构。改变允许的箭头，分类的结论就可能改变。

\section{把正交对角化读成一个同构}

本章使用有限维复内积空间；实数情形只需把“幺正”换成“正交”，把 $U^*$ 换成 $U^T$。称一个线性映射 $U$ 为幺正同构，意思是它是双射且保持内积。由此 $U^{-1}=U^*$，也有 $\|Ux-Uy\|=\|x-y\|$。

定义范畴 $\mathcal H$：对象是内积空间 $H$，箭头是幺正同构。定义范畴 $\mathcal S$：对象是带有自伴算子的内积空间 $(H,A)$，其中 $A^*=A$；从 $(H,A)$ 到 $(K,B)$ 的箭头 $U$ 除了要幺正，还必须满足
\[
 BU=UA.
\]
这个条件称为\textbf{交织}：先作用 $A$ 再换坐标，与先换坐标再作用 $B$，结果相同。恒等映射满足条件。若另有箭头 $V:(K,B)\to(L,C)$，则 $C(VU)=VBU=VUA$，所以复合仍是箭头；$BU=UA$ 也给出 $AU^{-1}=U^{-1}B$，所以逆映射仍是箭头。所有箭头都可逆的范畴叫\textbf{群胚}，因此 $\mathcal H,\mathcal S$ 都是群胚。

\begin{example}{同一个 $Q$ 为什么既保距又交织}
设 $A=\begin{pmatrix}5&2\\2&2\end{pmatrix}$，$D=\operatorname{diag}(6,1)$，$Q=\frac1{\sqrt5}\begin{pmatrix}2&-1\\1&2\end{pmatrix}$。证明 $Q$ 是从 $(\mathbb C^2,D)$ 到 $(\mathbb C^2,A)$ 的箭头，并解释在实平面中的几何意义。
\end{example}
\begin{solution}
$Q$ 的两列都是单位向量，内积为 $(-2+2)/5=0$，因此 $Q^*Q=I$。第一列乘 $A$ 变成自身的 $6$ 倍，第二列乘 $A$ 不变，于是 $AQ=QD$。交换图为
\[
\begin{array}{ccc}
 \mathbb C^2 &\xrightarrow{D}& \mathbb C^2\\
 \big\downarrow\rlap{$\scriptstyle Q$}&&\big\downarrow\rlap{$\scriptstyle Q$}\\
 \mathbb C^2 &\xrightarrow{A}& \mathbb C^2.
\end{array}
\]
两条路径给出的等式是 $AQ=QD$；$Q^*$ 是反向箭头。由于 $Q$ 的元素是实数，它和它的逆都把实向量送到实向量。限制到实平面，令 $x=Qu$，就有 $x^TAx=u^TQ^TAQu=u^TDu$，同时任意点对距离不变。故 $6u_1^2+u_2^2=1$ 的椭圆被保距地送到 $x^TAx=1$ 的椭圆。一般的复幺正矩阵未必保持实平面，例如 $iI$ 会把非零实向量送成纯虚向量。因此，从复空间的箭头回到实几何时，要确认这个实子空间确实被保留。
\end{solution}

\begin{example}{同构分类要记录特征值的重数}
在 $\mathbb C^3$ 中，比较 $D_1=\operatorname{diag}(6,6,1)$、$D_2=\operatorname{diag}(1,6,6)$、$D_3=\operatorname{diag}(6,1,1)$。确定哪些带算子的对象同构，并说明一般判据。
\end{example}
\begin{solution}
设 $e_1,e_2,e_3$ 为标准基。令 $Ue_1=e_2,\ Ue_2=e_3,\ Ue_3=e_1$，则 $U$ 把一组标准正交基送到另一组，故幺正。逐列检查得到 $D_2U=UD_1$，所以前两个对象同构。

若 $BU=UA$，则 $Av=\lambda v$ 推出 $B(Uv)=\lambda Uv$；对 $U^{-1}$ 再作同样的推理，$U$ 就把 $\ker(A-\lambda I)$ 一一映到 $\ker(B-\lambda I)$。$D_1$ 的 $6$ 特征子空间为二维，$D_3$ 的为一维，故两者不同构，尽管它们不计重数的特征值集合都是 $\{1,6\}$。

反之，若每个 $\lambda$ 的特征子空间维数相同，分别在对应的子空间中取标准正交基，并把每组基按顺序一一对应。由谱定理，这些基合起来就是全空间的标准正交基，因此得到的 $U$ 幺正；逐个基向量检查便有 $BU=UA$。所以在有限维自伴的条件下，分类数据是每个特征值及其重数。
\end{solution}

\section{忘却一部分结构与保留谱分量}

\begin{example}{忘掉算子后，允许的箭头变多了}
令 $F$ 将 $(H,A)$ 送到 $H$，将箭头 $U$ 送到同一个线性映射 $U$。证明 $F:\mathcal S\to\mathcal H$ 是函子。给出一个只在忘却后才允许的箭头。
\end{example}
\begin{solution}
一个\textbf{函子}要同时处理对象和箭头，并保持恒等箭头和复合。$\mathcal S$ 的每个箭头本来就是幺正同构，因此 $F(U)$ 的确是 $\mathcal H$ 的箭头；同一个恒等映射忘掉算子后仍是恒等映射。由于 $F$ 不改线性映射本身，$F(VU)=F(V)F(U)$。这就验证了函子条件。此处只忘掉附加的算子，故称为忘却函子。

取前例的 $(\mathbb C^2,D)$ 与 $(\mathbb C^2,A)$。两者经 $F$ 都变成 $\mathbb C^2$，而 $I$ 当然是 $\mathcal H$ 中的箭头。若这个 $I$ 来自原对象之间的箭头，就须有 $AI=ID$，即 $A=D$，矛盾。这说明在这对对象之间，忘却后的箭头多于原来箭头的像。术语上，$F$ \textbf{不满}；但原来两条不同的箭头仍是两个不同的线性映射，故 $F$ 对每对对象的箭头集合是单射，称为\textbf{忠实}。
\end{solution}

\begin{example}{谱投影与换坐标的两条路径}
仍取 $AQ=QD$。令 $P_6(A)$、$P_6(D)$ 分别为到 $6$ 特征子空间的正交投影。证明 $P_6(A)Q=QP_6(D)$，并用 $y=(1,2)^T$ 验算。
\end{example}
\begin{solution}
多项式 $p(t)=(t-1)/5$ 在 $6$ 处取 $1$，在 $1$ 处取 $0$。在特征基上计算可知 $P_6(A)=p(A)$、$P_6(D)=p(D)$，于是
\[
 P_6(A)=\frac15\begin{pmatrix}4&2\\2&1\end{pmatrix},\qquad
 P_6(D)=\begin{pmatrix}1&0\\0&0\end{pmatrix}.
\]
由 $AQ=QD$ 得 $(A-I)Q=Q(D-I)$，即所求等式。具体地，$Qy=(0,\sqrt5)^T$，所以 $P_6(A)Qy=(2,1)^T/\sqrt5$；另一条路径先得 $P_6(D)y=(1,0)^T$，再乘 $Q$ 也得 $(2,1)^T/\sqrt5$。

一般地，交织箭头把每个特征子空间幺正地送到同一特征值的子空间，并保持其正交补，故 $P_\lambda(B)U=UP_\lambda(A)$（没有该特征值时取零投影）。把每个对象 $(H,A)$ 忘却为向量空间 $H$，记这一函子为 $G$；它的目标范畴允许所有线性映射。上述交换条件恰好说 $P_\lambda$ 是从 $G$ 到 $G$ 的一个\textbf{自然变换}：每个空间上的分量是一个投影，且这些分量与所有允许的换坐标相容。投影往往不可逆，因此这里的目标范畴不能只允许幺正同构。
\end{solution}

\section{正交和与选择的不唯一性}

\begin{example}{正交直和如何合并重复特征值}
把 $A$ 与一维算子 $[6]$ 做正交直和。求 $A\oplus[6]$ 的特征子空间，并构造一个混合两块而仍与它交织的幺正变换。
\end{example}
\begin{solution}
正交直和的内积是两块内积之和。若两块各有交织箭头 $U,V$，则 $U\oplus V$ 保持这个内积，并且
\[
 (B\oplus B')(U\oplus V)=(U\oplus V)(A\oplus A')
\]
由两块各自的交织等式逐项得到。对本题，$Q\oplus[1]$ 给出对角形 $D\oplus[6]=\operatorname{diag}(6,1,6)$ 到 $A\oplus[6]$ 的箭头。设 $q_1,q_2$ 是 $Q$ 的两列，$v=(q_1,0)$、$w=(0,0,1)^T$、$s=(q_2,0)$。则 $6$ 的特征子空间由 $v,w$ 张成，$1$ 的由 $s$ 张成。特征值重数按块相加。

令 $Rv=(v+w)/\sqrt2$、$Rw=(-v+w)/\sqrt2$、$Rs=s$。三个像仍是标准正交基，所以 $R$ 幺正。算子 $A\oplus[6]$ 在 $v,w$ 所张成的空间上就是 $6I$，在 $s$ 上就是 $I$；分别在这三个基向量上比较两种作用顺序，可得 $R(A\oplus[6])=(A\oplus[6])R$。因此来自不同块的同一特征值分量可以互相混合。这里的“正交和”是对对象和箭头的逐块构造；通常的坐标嵌入并非幺正同构，不能直接把它当成这个群胚中的余积结构。
\end{solution}

\begin{example}{能否自然地选定一个单位特征向量}
假设要在每个非零的 $(H,A)$ 中选一个单位特征向量 $q_A$，并要求对任意交织箭头 $U:(H,A)\to(K,B)$，都有 $Uq_A=q_B$。这样的选择是否可能？
\end{example}
\begin{solution}
只看一维对象 $(\mathbb C,[6])$。它的每个单位复数都是特征向量。取 $U=-I$，则 $U$ 幺正且和 $[6]$ 交换，所以是此对象到自身的箭头。相容要求此时成为 $-q_A=q_A$，只能推出 $q_A=0$，和单位长度矛盾。这说明即使特征值只有一个、重数也是一，也不能按上述要求自然地挑选特征向量，更不能自然地指定一整组有序特征基。

另一方面，这个对象上的 $6$ 谱投影就是恒等映射，不因 $U=-I$ 改变。多维中也一样：把单位特征向量 $q$ 换成 $e^{i\theta}q$，投影 $qq^*$ 不变；有重数时，在同一子空间内更换标准正交基，也不改变到整个子空间的投影，因为投影由“在子空间上恒等、在正交补上为零”唯一确定。谱定理保证特征基\textbf{存在}，而真正与换坐标相容的是特征子空间及其谱投影。
\end{solution}
